\documentclass[12pt,a4paper]{article}

% ---- Packages ----
\usepackage[margin=1in]{geometry}
\usepackage{graphicx}
\usepackage{listings}
\usepackage{xcolor}
\usepackage{hyperref}
\usepackage{amsmath}
\usepackage{booktabs}
\usepackage{fancyhdr}
\usepackage{titlesec}
\usepackage{float}
\usepackage{caption}

% ---- Colors ----
\definecolor{codegreen}{rgb}{0,0.6,0}
\definecolor{codepurple}{rgb}{0.58,0,0.82}
\definecolor{codegray}{rgb}{0.5,0.5,0.5}
\definecolor{backcolour}{rgb}{0.95,0.95,0.95}
\definecolor{darkblue}{rgb}{0.0,0.0,0.6}

% ---- Code Listings Style ----
\lstdefinestyle{verilogstyle}{
    backgroundcolor=\color{backcolour},
    commentstyle=\color{codegreen},
    keywordstyle=\color{darkblue}\bfseries,
    numberstyle=\tiny\color{codegray},
    stringstyle=\color{codepurple},
    basicstyle=\ttfamily\scriptsize,
    breaklines=true,
    numbers=left,
    numbersep=5pt,
    frame=single,
    tabsize=4,
    captionpos=b,
    morekeywords={module, endmodule, input, output, wire, reg, always, posedge, negedge, assign, begin, end, if, else, parameter, localparam}
}
\lstset{style=verilogstyle}

% ---- Header / Footer ----
\pagestyle{fancy}
\fancyhf{}
\rhead{Assignment 9 -- Data Synchronizer}
\lhead{Digital IC Design}
\rfoot{Page \thepage}

% ---- Title Format ----
\titleformat{\section}{\Large\bfseries\color{darkblue}}{\thesection}{1em}{}
\titleformat{\subsection}{\large\bfseries\color{darkblue!80}}{\thesubsection}{1em}{}
\titleformat{\subsubsection}{\normalsize\bfseries\color{darkblue!60}}{\thesubsubsection}{1em}{}

% ---- Hyperref ----
\hypersetup{
    colorlinks=true,
    linkcolor=darkblue,
    urlcolor=blue,
    citecolor=darkblue
}

\begin{document}

% ===========================================================================
% TITLE PAGE
% ===========================================================================
\begin{titlepage}
\centering
\vspace*{2cm}
{\Huge\bfseries Data Synchronizer\\(DATA\_SYNC)\par}
\vspace{1cm}
{\Large Assignment 9 -- Digital IC Design Course\par}
\vspace{2cm}
{\large\textbf{Submitted by:} Ziad Mostafa Abdelaziz\par}
\vspace{0.5cm}
{\large\today\par}
\vfill
\end{titlepage}

\tableofcontents
\newpage

% ===========================================================================
% 1. INTRODUCTION
% ===========================================================================
\section{Introduction}

Clock Domain Crossing (CDC) in the digital domain is defined as the process of passing a signal or vector (multi-bit signal) from one clock domain to another clock domain, which introduces many issues such as \textbf{metastability}, \textbf{data incoherence}, and \textbf{data loss}.

This assignment implements the \textbf{Synchronized MUX-Select Synchronization Scheme}, which is used to synchronize multiple bits safely across clock domains. The scheme uses a multi-stage flip-flop synchronizer to safely transfer a single-bit enable signal, and a MUX-registered data path to atomically capture the multi-bit data bus when the synchronized enable pulse fires.

% ===========================================================================
% 2. BLOCK INTERFACE
% ===========================================================================
\section{Block Interface}

\subsection{Parameters}
\begin{table}[H]
\centering
\begin{tabular}{@{}ll@{}}
\toprule
\textbf{Parameter Name} & \textbf{Description} \\
\midrule
\texttt{NUM\_STAGES} & Number of Flip Flop Stages (default = 2) \\
\texttt{BUS\_WIDTH}  & Width of synchronized bus (default = 8) \\
\bottomrule
\end{tabular}
\caption{Design Parameters}
\end{table}

\subsection{Ports Description}
\begin{table}[H]
\centering
\begin{tabular}{@{}llc@{}}
\toprule
\textbf{Signal Name} & \textbf{Description} & \textbf{Width} \\
\midrule
\texttt{unsync\_bus}   & Unsynchronized bus from source domain & Parameterized (default = 8) \\
\texttt{bus\_enable}   & Source domain enable signal            & 1 \\
\texttt{CLK}           & Destination domain clock               & 1 \\
\texttt{RST}           & Destination domain active-low async reset & 1 \\
\texttt{sync\_bus}     & Synchronized bit/bus                  & Parameterized (default = 8) \\
\texttt{enable\_pulse} & Destination domain enable signal      & 1 \\
\bottomrule
\end{tabular}
\caption{Port Descriptions}
\end{table}

% ===========================================================================
% 3. ARCHITECTURE
% ===========================================================================
\section{Architecture and Design}

The DATA\_SYNC module implements the Synchronized MUX-Select scheme with four main components:

\begin{enumerate}
    \item \textbf{Multi Flip-Flop Synchronizer Chain:} A parameterized chain of \texttt{NUM\_STAGES} flip-flops synchronizes the \texttt{bus\_enable} signal from the source clock domain into the destination clock domain, mitigating metastability. The chain is implemented as a shift register:
    \[
        \texttt{sync\_stage\_reg} \leftarrow \{\texttt{sync\_stage\_reg[NUM\_STAGES-2:0]},\; \texttt{bus\_enable}\}
    \]
    
    \item \textbf{Pulse Generator (Rising-Edge Detector):} Compares the synchronized enable output with a one-cycle delayed copy to produce a single-cycle \texttt{generated\_pulse}:
    \[
        \texttt{generated\_pulse} = \texttt{sync\_stage\_reg[NUM\_STAGES-1]} \;\land\; \overline{\texttt{enable\_flop}}
    \]
    This ensures the pulse fires exactly once on the $0 \to 1$ transition of the synchronized enable.
    
    \item \textbf{MUX-Select Registered Data Bus:} A 2-to-1 multiplexer controlled by \texttt{generated\_pulse} selects between the current \texttt{sync\_bus} value (hold) and \texttt{unsync\_bus} (update). The output is registered, ensuring \texttt{sync\_bus} updates atomically.
    
    \item \textbf{Enable Pulse Register:} Registers \texttt{generated\_pulse} to produce the output \texttt{enable\_pulse}, which is synchronized with the \texttt{sync\_bus} data update on the same clock edge.
\end{enumerate}

% ===========================================================================
% 4. RTL CODE
% ===========================================================================
\section{RTL Verilog Code}

\lstinputlisting[language=Verilog, caption={DATA\_SYNC.v -- Data Synchronizer RTL Module}]{../DATA_SYNC.v}

% ===========================================================================
% 5. TESTBENCH
% ===========================================================================
\section{Testbench}

The self-checking testbench verifies the DATA\_SYNC module under an asynchronous dual-clock setup:
\begin{itemize}
    \item \textbf{Source Clock:} $f_{tx} = 50$ MHz (Period = 20 ns)
    \item \textbf{Destination Clock:} $f_{dest} = 100$ MHz (Period = 10 ns)
\end{itemize}

A concurrent \texttt{always} block monitor captures \texttt{enable\_pulse} events and the corresponding \texttt{sync\_bus} data in real-time, avoiding race conditions between the source and destination clock domains.

\subsection{Test Cases}
\begin{table}[H]
\centering
\begin{tabular}{@{}clc@{}}
\toprule
\textbf{TC\#} & \textbf{Description} & \textbf{Expected Result} \\
\midrule
1 & Single Transfer (\texttt{0xA5})       & \texttt{sync\_bus} = \texttt{0xA5}, 1-cycle pulse \\
2 & Consecutive (\texttt{0x5A} $\to$ \texttt{0xF0}) & Both values synchronized correctly \\
3 & Dynamic Transfer (\texttt{0xC3})      & \texttt{sync\_bus} = \texttt{0xC3}, 1-cycle pulse \\
4 & All-Ones (\texttt{0xFF})              & \texttt{sync\_bus} = \texttt{0xFF}, 1-cycle pulse \\
5 & All-Zeros (\texttt{0x00})             & \texttt{sync\_bus} = \texttt{0x00}, 1-cycle pulse \\
6 & Async Active-Low Reset               & Immediate clear of all outputs \\
\bottomrule
\end{tabular}
\caption{Test Case Summary}
\end{table}

\lstinputlisting[language=Verilog, caption={DATA\_SYNC\_tb.v -- Self-Checking Testbench}]{../DATA_SYNC_tb.v}

% ===========================================================================
% 6. SIMULATION RESULTS
% ===========================================================================
\section{Simulation Results}

\subsection{ModelSim Transcript Log}

All 6 test cases passed successfully with 0 errors. The simulation completed in 1285 ns.

\begin{figure}[H]
    \centering
    \includegraphics[width=\textwidth]{log output.png}
    \caption{ModelSim Transcript Log -- All 6 Test Cases PASSED (0 Errors)}
\end{figure}

\subsection{Output Waveform}

\begin{figure}[H]
    \centering
    \includegraphics[width=\textwidth]{output waveform.png}
    \caption{ModelSim Output Waveform -- DATA\_SYNC Simulation}
\end{figure}

\subsection{Detailed Waveform and Signal Transition Analysis}

The output waveform captures the complete simulation of all six test cases. Each signal group is analyzed below:

\subsubsection{Source Domain Signals}

\begin{itemize}
    \item \textbf{\texttt{tx\_clk} (Cyan):} The source clock runs at $f_{tx} = 50$ MHz with a period of $T_{tx} = 20$ ns. This clock drives the \texttt{unsync\_bus} and \texttt{bus\_enable} signals.
    
    \item \textbf{\texttt{bus\_enable} (Yellow):} The source-domain enable signal pulses high for 4 source-clock cycles ($4 \times 20 = 80$ ns) per data transfer, then returns to 0. The waveform shows 7 distinct \texttt{bus\_enable} pulses corresponding to TC1 through TC6 (TC2 has two consecutive transfers). Between pulses, there is a sufficient idle gap for the synchronizer chain to fully propagate the falling edge before the next rising edge.
    
    \item \textbf{\texttt{unsync\_bus} (White, Hex):} The multi-bit data bus changes values at the source clock domain: \texttt{0xA5} $\to$ \texttt{0x5A} $\to$ \texttt{0xF0} $\to$ \texttt{0xC3} $\to$ \texttt{0xFF} $\to$ \texttt{0x00} $\to$ \texttt{0xBB} (reset test). Each value is held stable by the testbench for the entire duration of the \texttt{bus\_enable} pulse, ensuring the data is valid when the destination domain MUX captures it.
\end{itemize}

\subsubsection{Destination Domain Control Signals}

\begin{itemize}
    \item \textbf{\texttt{CLK} (Cyan):} The destination domain clock runs at $f_{dest} = 100$ MHz with a period of $T_{dest} = 10$ ns -- twice the frequency of the source clock. All internal synchronizer registers and output registers are clocked by this signal.
    
    \item \textbf{\texttt{RST} (Orange):} Active-low asynchronous reset. Stays high during normal operation. In TC6 (Async Reset test), \texttt{RST} drops to 0 mid-operation, causing an immediate asynchronous clear of all internal registers (\texttt{sync\_stage\_reg}, \texttt{enable\_flop}, \texttt{sync\_bus}, \texttt{enable\_pulse}) regardless of the clock state.
\end{itemize}

\subsubsection{Internal Synchronizer Signals}

\begin{itemize}
    \item \textbf{\texttt{sync\_stage\_reg} (Magenta, Binary):} The 2-stage flip-flop synchronizer shift register. For each \texttt{bus\_enable} pulse:
    \begin{enumerate}
        \item Initially \texttt{2'b00} (idle).
        \item After 1 destination clock edge: \texttt{2'b01} (bit 0 samples \texttt{bus\_enable} = 1).
        \item After 2 destination clock edges: \texttt{2'b11} (bit 1 propagates the synchronized enable).
        \item After \texttt{bus\_enable} deasserts: returns to \texttt{2'b10} $\to$ \texttt{2'b00} over 2 destination clocks.
    \end{enumerate}
    The waveform clearly shows this \texttt{00} $\to$ \texttt{11} $\to$ \texttt{00} transition pattern repeating for each data transfer.
    
    \item \textbf{\texttt{enable\_flop} (Yellow):} A 1-cycle delayed copy of \texttt{sync\_stage\_reg[1]}. Rises one destination clock cycle \emph{after} \texttt{sync\_stage\_reg} reaches \texttt{2'b11}. Falls one cycle after \texttt{sync\_stage\_reg[1]} drops back to 0.
    
    \item \textbf{\texttt{generated\_pulse} (Purple):} The combinational rising-edge pulse:
    \[
        \texttt{generated\_pulse} = \texttt{sync\_stage\_reg[1]} \;\land\; \overline{\texttt{enable\_flop}}
    \]
    The waveform shows \texttt{generated\_pulse} as a clean single-cycle spike exactly on the clock cycle where \texttt{sync\_stage\_reg[1]} first becomes 1 while \texttt{enable\_flop} is still 0. This pulse is exactly 1 destination clock period wide ($T_{dest} = 10$ ns).
\end{itemize}

\subsubsection{Destination Domain Output Signals}

\begin{itemize}
    \item \textbf{\texttt{sync\_bus} (Green, Hex):} The synchronized multi-bit output bus. Initially \texttt{0x00} after reset. The waveform shows \texttt{sync\_bus} updating to each new data value exactly when \texttt{generated\_pulse} fires:
    \begin{itemize}
        \item \texttt{0x00} $\to$ \texttt{0xA5} (TC1): Data captured correctly.
        \item \texttt{0xA5} $\to$ \texttt{0x5A} $\to$ \texttt{0xF0} (TC2): Both consecutive values captured in sequence.
        \item \texttt{0xF0} $\to$ \texttt{0xC3} (TC3): Dynamic transfer verified.
        \item \texttt{0xC3} $\to$ \texttt{0xFF} (TC4): All-ones captured correctly.
        \item \texttt{0xFF} $\to$ \texttt{0x00} (TC5): All-zeros captured correctly.
        \item \texttt{0x00} $\to$ \texttt{0xBB} $\to$ \texttt{0x00} (TC6): Data loaded, then cleared by async reset.
    \end{itemize}
    Between \texttt{generated\_pulse} events, \texttt{sync\_bus} holds its previous value (the MUX selects the feedback path), ensuring data stability.
    
    \item \textbf{\texttt{enable\_pulse} (Lime):} A registered version of \texttt{generated\_pulse}. The waveform confirms that \texttt{enable\_pulse} is a single destination-clock-cycle wide pulse that fires exactly one clock cycle after \texttt{generated\_pulse}. Since both \texttt{sync\_bus} and \texttt{enable\_pulse} are clocked by the same edge, they are guaranteed to be valid simultaneously -- the downstream logic can safely use \texttt{enable\_pulse} as a write-enable for \texttt{sync\_bus}.
\end{itemize}

\subsubsection{Synchronization Latency Analysis}

The total synchronization latency from \texttt{bus\_enable} assertion (source domain) to \texttt{enable\_pulse} assertion (destination domain) is:

\[
    t_{sync} = (\texttt{NUM\_STAGES} + 2) \times T_{dest} + t_{setup}
\]

For the default configuration (\texttt{NUM\_STAGES} = 2, $T_{dest} = 10$ ns):
\begin{itemize}
    \item 2 cycles for the synchronizer chain to propagate \texttt{bus\_enable} through both flip-flops.
    \item 1 cycle for \texttt{enable\_flop} delay (pulse generator).
    \item 1 cycle for registering \texttt{generated\_pulse} into \texttt{enable\_pulse}.
    \item Total: $\approx 4 \times 10 = 40$ ns (plus metastability resolution time $t_{setup}$).
\end{itemize}

\subsubsection{Asynchronous Reset Behavior (TC6)}

In TC6, the waveform shows:
\begin{itemize}
    \item \texttt{sync\_bus} = \texttt{0xBB} is loaded via a normal transfer.
    \item \texttt{RST} drops to 0 asynchronously (between clock edges).
    \item Within less than 1 ns, all registers clear: \texttt{sync\_stage\_reg} $\to$ \texttt{2'b00}, \texttt{enable\_flop} $\to$ 0, \texttt{sync\_bus} $\to$ \texttt{0x00}, \texttt{enable\_pulse} $\to$ 0.
    \item After \texttt{RST} returns to 1, the module resumes normal operation.
\end{itemize}

% ===========================================================================
% 7. SPYGLASS LINT RESULTS
% ===========================================================================
\section{SpyGlass Lint Results}

The design was analyzed using Synopsys SpyGlass (version vL-2016.06) with the \texttt{rtl\_handoff} methodology. The lint analysis covers multiple policies: lint, morelint, clock-reset, ERC, latch, simulation, STARC, STARC2005, timing, and TXV.

\subsection{Design Read Step Output}

SpyGlass successfully read and elaborated the \texttt{DATA\_SYNC.v} source file. The message tree shows \textbf{Total: 2 messages, Displayed: 2, Waived: 0} -- both are informational \texttt{Design Read} messages (no errors or warnings).

\begin{figure}[H]
    \centering
    \includegraphics[width=\textwidth]{run design read step output from spyglass.png}
    \caption{SpyGlass Design Read Step Output -- 2 Info Messages, 0 Errors}
\end{figure}

\subsection{Lint RTL Goal Run Output}

The Lint RTL goal run completed with a \textbf{clean lint result}. The message tree again shows \textbf{Total: 2 messages, Displayed: 2, Waived: 0} -- both are informational messages from the Design Read phase, confirming:

\begin{itemize}
    \item \textbf{0 Errors}: No design rule violations.
    \item \textbf{0 Warnings}: No potential issues detected.
    \item \textbf{0 Latches}: No unintended latches inferred.
\end{itemize}

\begin{figure}[H]
    \centering
    \includegraphics[width=\textwidth]{lint rtl goal run output clean lint spyglass.png}
    \caption{SpyGlass Lint RTL Goal Run Output -- Clean Lint (0 Errors, 0 Warnings)}
\end{figure}

\subsection{SpyGlass Lint Summary Report (moresimple.rpt)}

\begin{lstlisting}[language={}, caption={SpyGlass moresimple.rpt -- Lint Summary Report}]
Total Number of Generated Messages :          2
Number of Waived Messages          :          0
Number of Reported Messages        :          2
Number of Overlimit Messages       :          0

MORESIMPLE REPORT:

############### BuiltIn -> RuleGroup=Design Read ###############
[1] DetectTopDesignUnits  Info  DATA_SYNC.v  Line 15
    Module DATA_SYNC is a top level design unit
[0] ElabSummary           Info
    Please refer file elab_summary.rpt for elab summary report
\end{lstlisting}

Both reported messages are \textbf{Info} severity -- purely informational notifications confirming the top-level module detection and elaboration summary. No errors, warnings, or latch inferences were detected across all lint policies.

% ===========================================================================
% 8. CONCLUSION
% ===========================================================================
\section{Conclusion}

The DATA\_SYNC module successfully implements the Synchronized MUX-Select Synchronization Scheme for safe multi-bit Clock Domain Crossing. Key achievements:

\begin{itemize}
    \item \textbf{Parameterized design} supporting configurable synchronizer stages (\texttt{NUM\_STAGES}) and bus width (\texttt{BUS\_WIDTH}).
    \item \textbf{Rising-edge pulse generator} ensures a single-cycle enable pulse for precise data capture timing.
    \item \textbf{MUX-registered data path} ensures \texttt{sync\_bus} updates atomically with \texttt{enable\_pulse}.
    \item \textbf{Active-low asynchronous reset} clears all internal registers immediately.
    \item \textbf{All 6 test cases PASSED} with 0 errors in ModelSim simulation.
    \item \textbf{Clean SpyGlass lint}: 0 errors, 0 warnings, 0 latches across all lint policies.
\end{itemize}

\end{document}
